\begin{figure}[tbp]
\centering
\begin{tikzpicture}[x=1cm,y=1cm]
\node[box,text width=2.3cm] (dev) at (0,0) {串口模型\\接收 / 状态更新};
\node[box,text width=2.3cm] (plic) at (3.3,0) {PLIC 模型\\pending / 优先级};
\node[box,text width=2.3cm] (pending) at (6.6,0) {hart 中断输入\\CPU pending 状态};
\node[amberbox,text width=2.3cm] (kick) at (9.9,0) {通知执行线程\\退出请求 / kick};
\draw[wire] (dev)--(plic); \draw[wire] (plic)--(pending); \draw[wire] (pending)--(kick);
\node[box,text width=2.9cm] (timer) at (6.6,1.65) {CLINT 定时器 / 软件中断};
\draw[wire] (timer)--(pending);
\node[amberbox,text width=2.3cm] (check) at (9.9,-2.0) {按架构检查\\屏蔽 / 委托 / 特权};
\node[greenbox,text width=2.3cm] (state) at (6.6,-2.0) {接收中断\\更新 CSR 与 PC};
\node[greenbox,text width=2.3cm] (isr) at (3.3,-2.0) {执行 guest\\中断处理程序};
\node[box,text width=2.3cm] (ack) at (0,-2.0) {PLIC claim / complete\\及设备应答};
\draw[wire] (kick)--(check); \draw[wire] (check)--node[above,font=\footnotesize]{可接收}(state);
\draw[wire] (state)--(isr); \draw[wire] (isr)--(ack);
\draw[dwire] (ack.north)--node[midway,above,sloped,font=\footnotesize]{更新控制器状态}(plic.south);
\end{tikzpicture}
\caption[QEMU 中断生成与投递]{QEMU 中断生成与投递。外部中断以 UART/PLIC 为例；本地定时器和软件中断不经过 PLIC。置 pending、通知线程和 guest 接收中断是不同动作。\cite{source3}}
\label{fig:softirq}
\end{figure}
